\daytitle{15}{电荷、电场与电场强度}

\begin{learningbox}{本节导读}
电磁学的第一块基石。从电荷、库仑定律开始，引入场的概念——电场、电场强度、电场线。理解场是物质存在的一种形式，以及如何用“场”的观点代替“超距作用”。
\end{learningbox}

\textbf{本节读题路线：} 先分清源电荷和试探电荷，再判断要求的是电场力还是电场强度；点电荷问题先画方向，再按矢量叠加求合场强。遇到电场线，先用箭头读方向、用疏密读强弱；遇到静电平衡题，再检查导体内部、表面和空腔中的场强特点。

\section*{一、学习目标}

\begin{enumerate}
  \item 能用电荷守恒解释常见起电方式，并使用库仑定律 $F = k\dfrac{|Q_1 Q_2|}{r^2}$。
  \item 能用 $\vect E = \dfrac{\vect F}{q_0}$ 定义电场强度，并用电场线描述电场方向和强弱。
  \item 能由库仑定律和场强定义得到点电荷场强 $E = k\dfrac{|Q|}{r^2}$，并能判断方向。
  \item 能用矢量叠加原理分析多个点电荷产生的合场强。
  \item 区分匀强电场和非匀强电场的特点。
  \item 能用电场线和电荷分布说明导体静电平衡、静电屏蔽的基本结论。
  \item 会把电场力与重力等恒力合成，初步建立“等效重力场”模型。
\end{enumerate}
\section*{二、初高中衔接}

初中物理讲过“电荷”和“电荷间的相互作用”——同种电荷相斥、异种电荷相吸。这是定性的。

高中在三个方向上深化：
\begin{enumerate}
  \item \textbf{定量化：} 库仑定律给出了电荷间作用力的精确公式。
  \item \textbf{场的概念：} 引入“电场”——电荷周围存在的一种看不见的物质，通过电场传递作用力。
  \item \textbf{场的强弱：} 电场强度 $E$ 定量描述电场“有多强”和“朝哪指”。
\end{enumerate}

\textbf{科学发展视角：} 电场模型贯穿现代电子技术：触摸屏靠电场分布变化感知手指位置，静电除尘利用电场控制带电微粒，芯片内部的电子运动也离不开电场控制。高中先学点电荷和匀强电场，是为了建立可计算、可预测的场模型。

\section*{三、核心概念}

\begin{conceptbox}{电荷、起电与电荷守恒}

\textbf{电荷：}
\begin{itemize}
  \item 正电荷（$+$）和负电荷（$-$），同种相斥、异种相吸。
  \item 电荷量 $Q$ 的单位：库仑（$\mathrm{C}$）。
  \item 元电荷 $e = 1.60 \times 10^{-19}\ \mathrm{C}$。
  \item 电荷量具有量子化特征：带电体的净电荷量是元电荷的整数倍，即 $Q=ne$（$n$ 为整数）。
  \item 电荷守恒定律：孤立系统的总电荷量保持不变。高中常见起电现象来自电荷转移或重新分布，而不是凭空产生净电荷。
\end{itemize}

\textbf{常见起电方式：}
\begin{itemize}
  \item \textbf{摩擦起电：} 两种物体摩擦时，电子从一个物体转移到另一个物体，两个物体常带异种电荷。
  \item \textbf{接触起电：} 带电体接触导体后，电荷在接触的导体间重新分配。
  \item \textbf{感应起电：} 带电体靠近导体时，导体内部自由电荷重新分布；若配合接地或分开导体，可使导体带电。
\end{itemize}

\begin{center}
\begin{tikzpicture}[x=1cm,y=1cm,>=Stealth]
  \begin{scope}
    \node[font=\small\bfseries] at (-3.65,1.45) {摩擦起电};
    \draw[rounded corners=4pt,thick,fill=blue!8] (-4.85,-0.25) rectangle (-3.75,0.35);
    \draw[rounded corners=4pt,thick,fill=orange!12] (-3.45,-0.25) rectangle (-2.35,0.35);
    \draw[->,thick] (-4.05,0.62) -- (-3.15,0.62);
    \draw[->,thick] (-3.15,-0.52) -- (-4.05,-0.52);
    \foreach \x in {-4.55,-4.25,-3.95}{\node[blue] at (\x,0.05) {$-$};}
    \foreach \x in {-3.15,-2.85,-2.55}{\node[red] at (\x,0.05) {$+$};}
    \node[font=\scriptsize,align=center] at (-3.60,-1.05) {电子在摩擦中转移\\两物体带异种电荷};
  \end{scope}

  \begin{scope}
    \node[font=\small\bfseries] at (0,1.45) {接触起电};
    \draw[thick,fill=blue!8] (-1.05,0.05) circle (0.38);
    \draw[thick,fill=gray!10] (0.05,0.05) circle (0.38);
    \node[blue] at (-1.05,0.05) {$-$};
    \draw[->,thick] (-0.58,0.05) -- (-0.32,0.05);
    \draw[thick,fill=blue!8] (0.90,0.05) circle (0.38);
    \draw[thick,fill=blue!8] (1.72,0.05) circle (0.38);
    \node[blue] at (0.90,0.05) {$-$};
    \node[blue] at (1.72,0.05) {$-$};
    \node[font=\scriptsize,align=center] at (0.35,-1.05) {带电体接触导体\\电荷重新分配};
  \end{scope}

  \begin{scope}
    \node[font=\small\bfseries] at (4.05,1.45) {感应起电};
    \draw[thick,fill=red!8] (2.70,0.05) circle (0.34);
    \node[red] at (2.70,0.05) {$+$};
    \draw[thick,rounded corners=10pt,fill=gray!10] (3.45,-0.28) rectangle (5.05,0.38);
    \node[blue] at (3.75,0.06) {$-$};
    \node[blue] at (4.08,0.06) {$-$};
    \node[red] at (4.58,0.06) {$+$};
    \node[red] at (4.86,0.06) {$+$};
    \draw[->,thick] (4.95,-0.50) -- (4.95,-0.95);
    \node[font=\scriptsize] at (5.50,-0.75) {接地或分开};
    \node[font=\scriptsize,align=center] at (4.05,-1.35) {不接触也能使导体\\电荷重新分布};
  \end{scope}

  \node[font=\scriptsize,align=center] at (0,-2.15)
    {三种起电都不是创造电荷，而是电子转移或电荷重新分布，系统总电荷守恒};
\end{tikzpicture}


\figcaption{三种起电方式：本质是电子转移或电荷重新分布}\label{fig:charging-methods}
\end{center}

\textbf{认知提醒：} 起电不是“产生了新的电荷”。例如摩擦起电中，一个物体失去电子带正电，另一个物体得到电子带负电；若把两物体看成一个系统，总电荷仍守恒。
\end{conceptbox}

\begin{conceptbox}{库仑定律}
\textbf{库仑定律：}
\[
\boxed{F = k\frac{|Q_1 Q_2|}{r^2}}
\]

\begin{itemize}
  \item $k = 9.0 \times 10^9\ \mathrm{N \cdot m^2 / C^2}$（静电力常量）。
  \item $F$：静电斥力（$Q_1$、$Q_2$ 同号）或吸力（异号）。
  \item 适用条件：真空中的点电荷（高中可近似用于空气中的点电荷）。
\end{itemize}

\begin{center}
% TikZ 图：库仑定律（两个点电荷）
% Day 15：电荷、电场与电场强度
\begin{tikzpicture}[scale=0.8,>=Latex]
  % 两个正电荷
  \filldraw[red] (0,0) circle (4pt);
  \node[above] at (0,0) {$+Q_1$};
  \filldraw[red] (4,0) circle (4pt);
  \node[above] at (4,0) {$+Q_2$};
  % 距离
  \draw[<->] (0,-0.65) -- (4,-0.65);
  \node[below] at (2,-0.65) {$r$};
  % 斥力：同号电荷相互远离
  \draw[->,thick,blue] (-0.2,0) -- (-1.4,0) node[midway,above] {$\vect{F}$};
  \draw[->,thick,blue] (4.2,0) -- (5.4,0) node[midway,above] {$\vect{F}$};
  % 公式
  \node at (2,-1.75) {$F = k\dfrac{|Q_1 Q_2|}{r^2}$};
\end{tikzpicture}


\figcaption{库仑定律：同号电荷相斥，大小 $F = k\dfrac{|Q_1 Q_2|}{r^2}$}\label{fig:coulomb-law}
\end{center}

\end{conceptbox}

\begin{conceptbox}{电场与电场强度}

\textbf{电场：} 电荷周围存在的一种特殊物质。电场对放入其中的电荷有力的作用。

\textbf{电场强度（定义式）：}
\[
\boxed{\vect E = \frac{\vect F}{q_0}}\quad(q_0>0)
\]

\begin{itemize}
  \item 矢量，方向与正电荷在该点受到的电场力方向相同；其大小为 $E=\dfrac{F}{q_0}$。
  \item 单位：$\mathrm{N/C}$ 或 $\mathrm{V/m}$。
  \item 由电场本身决定，与正试探电荷 $q_0$ 无关。
  \item 试探电荷应足够小，放入后不能明显改变原来的电场分布。
\end{itemize}

\textbf{点电荷的电场强度：}
\[
E = k\frac{|Q|}{r^2}
\]

方向由场源电荷 $Q$ 的正负决定：正点电荷的电场向外，负点电荷的电场向内。

\textbf{公式来源与边界：} 对点电荷 $Q$，正试探电荷 $q_0$ 受到 $F=k\dfrac{|Q|q_0}{r^2}$，代入定义式 $E=\dfrac{F}{q_0}$ 才得到 $E=k\dfrac{|Q|}{r^2}$。前者是任意电场的定义方法，后者只是点电荷电场的计算结果，不能直接套到带电板、导体表面或复杂电荷分布上。

\end{conceptbox}

\begin{conceptbox}{电场叠加原理}

多个点电荷在某点产生的合场强等于各点电荷单独在该点产生的场强的矢量和。

\begin{center}
% TikZ 图：电场叠加
% Day 15：电荷、电场与电场强度
\begin{tikzpicture}[scale=0.9,>=Latex]
  % +Q 和 -Q
  \filldraw[red] (0,0) circle (3pt);
  \node[above] at (0,0) {$+Q$};
  \filldraw[blue] (3,0) circle (3pt);
  \node[above] at (3,0) {$-Q$};
  % 中点 P
  \filldraw[black] (1.5,0) circle (2pt);
  \node[below] at (1.5,0) {$P$};
  % E1（+Q 对 P 向右）
  \draw[->,thick,red] (1.5,0) -- (2.5,0) node[above] {$\vect{E}_1$};
  % E2（-Q 对 P 向右）
  \draw[->,thick,blue] (1.5,0) -- (2.5,-0.3) node[below] {$\vect{E}_2$};
  % 合电场
  \draw[->,ultra thick,black] (1.5,0) -- (3.2,0) node[above] {$\vect{E}$};
\end{tikzpicture}


\figcaption{电场叠加：$\vect{E} = \vect{E}_1 + \vect{E}_2$}\label{fig:electric-superposition}
\end{center}

\end{conceptbox}

\begin{conceptbox}{电场线}

电场线是描述电场分布的假想曲线：
\begin{itemize}
  \item 电场线从正电荷或无穷远出发，终止于负电荷或无穷远。
  \item 在同一幅按统一规则绘制的示意图中，电场线越密表示电场越强；这种疏密只作定性比较。
  \item 电场线某点的切线方向为该点的场强方向。
  \item 电场线不相交。
\end{itemize}

\begin{center}
% TikZ 图：点电荷的电场线
% Day 15：电荷、电场与电场强度
\begin{tikzpicture}[scale=0.8,>=Latex,line cap=round]
  % 正点电荷
  \filldraw[red] (0,0) circle (4pt);
  \node[below=5pt,fill=white,inner sep=1pt] at (0,0) {$+Q$};
  % 电场线（从正电荷向外辐射）
  \foreach \angle in {0,30,60,90,120,150,180,210,240,270,300,330} {
    \draw[->,thick] (\angle:0.18) -- (\angle:2.65);
  }
\end{tikzpicture}


\figcaption{正点电荷的电场线：从正电荷向外辐射}\label{fig:electric-field-lines}

\vspace{0.5em}
% TikZ 图：匀强电场
% Day 15：电荷、电场与电场强度
\begin{tikzpicture}[scale=0.9,>=Latex]
  % 平行板
  \draw[thick] (-2,1.5) -- (2,1.5);
  \node[above] at (0,1.5) {$+$};
  \draw[thick] (-2,-1.5) -- (2,-1.5);
  \node[below] at (0,-1.5) {$-$};
  % 电场线（竖直向下，均匀等距）
  \foreach \x in {-1.5,-0.5,0.5,1.5} {
    \draw[->,thick,blue] (\x,1.3) -- (\x,-1.3);
  }
  \node[blue,right] at (2.2,0) {$\vect{E}$};
\end{tikzpicture}


\figcaption{匀强电场：平行等距的电场线}\label{fig:uniform-field}
\end{center}

\end{conceptbox}

\begin{samepage}
\section*{四、公式与适用条件}

\textbf{公式使用前检查：}
\nopagebreak[4]
\begin{enumerate}
  \item 先判断场源电荷和试探电荷，不要把“产生电场的电荷”和“受力的电荷”混在一起。
  \item 电场强度是矢量，多个点电荷电场要按方向叠加，不能只加大小。
  \item 正电荷受力方向与电场方向相同，负电荷受力方向与电场方向相反。
  \item 画电场线时只表达电场方向和强弱趋势，不能把电场线当成真实轨迹。
\end{enumerate}
\end{samepage}

\begin{center}
\begin{tabular}{lll}
\toprule
公式 & 物理意义 & 适用条件 \\
\midrule
$Q=ne$ & 电荷量子化 & $n$ 为整数，$e$ 为元电荷 \\
$F = k\dfrac{|Q_1 Q_2|}{r^2}$ & 库仑力大小 & 真空点电荷 \\
$\vect E = \dfrac{\vect F}{q_0}$ & 场强定义 & $q_0>0$；大小为 $E=\dfrac{F}{q_0}$ \\
$\vect F=q\vect E$ & 电场力 & 大小为 $F=|q|E$，方向看电荷正负 \\
$E = k\dfrac{|Q|}{r^2}$ & 点电荷场强大小 & 点电荷周围 \\
$\vect{E} = \sum \vect{E}_i$ & 电场叠加 & 多个场源共同作用时的矢量和 \\
\bottomrule
\end{tabular}
\end{center}

\begin{learningbox}{电场题先分清三个角色}
电场题最容易把“产生场”“定义场”“在场中受力”混在一起。列式前先分清：
\begin{center}
\begin{tabular}{p{0.18\textwidth}p{0.28\textwidth}p{0.42\textwidth}}
\toprule
角色 & 作用 & 解题时看什么 \\
\midrule
场源电荷 $Q$ & 产生电场 & 点电荷场强 $E=k\dfrac{|Q|}{r^2}$ 由 $Q$ 和位置决定，方向由 $Q$ 的正负决定。 \\
正试探电荷 $q_0$ & 定义场强方向和大小 & $\vect E=\dfrac{\vect F}{q_0}$ 中约定 $q_0>0$，且足够小；它只是定义工具，不决定原电场。 \\
实际带电体 $q$ & 在电场中受力或运动 & 电场力 $\vect{F}=q\vect{E}$；$q>0$ 时受力沿电场方向，$q<0$ 时受力反向。 \\
\bottomrule
\end{tabular}
\end{center}
\textbf{读题提醒：} 求场强时先只看场源；求受力时再放入实际电荷；画电场线时画的是场的方向和强弱趋势，不是电荷轨迹。
\end{learningbox}

\section*{五、典型模型}

\subsection*{模型 1：等量异种电荷的电场}

中垂线上场强的方向平行于两电荷连线。中点处场强不为零（两电荷在该点产生的场强方向相同）。

\begin{center}
% TikZ 图：等量异种电荷的电场线
% Day 15：电荷、电场与电场强度
\begin{tikzpicture}[scale=0.9,>=Latex]
  \coordinate (P) at (-1.6,0);
  \coordinate (N) at (1.6,0);

  \filldraw[red] (P) circle (4pt);
  \node[below] at (P) {$+Q$};
  \filldraw[blue] (N) circle (4pt);
  \node[below] at (N) {$-Q$};

  \draw[->,thick] (-1.25,0) -- (1.25,0);
  \draw[->,thick] (-1.35,0.18) .. controls (-0.75,0.95) and (0.75,0.95) .. (1.35,0.18);
  \draw[->,thick] (-1.35,-0.18) .. controls (-0.75,-0.95) and (0.75,-0.95) .. (1.35,-0.18);
  \draw[->,thick] (-1.52,0.25) .. controls (-1.05,1.55) and (1.05,1.55) .. (1.52,0.25);
  \draw[->,thick] (-1.52,-0.25) .. controls (-1.05,-1.55) and (1.05,-1.55) .. (1.52,-0.25);

  \draw[dashed,gray] (0,-1.45) -- (0,1.45);
  \node[gray,above] at (0,1.45) {中垂线};
  \node[font=\small] at (0,-1.75) {中垂线上合场强平行于两电荷连线};
\end{tikzpicture}


\figcaption{等量异种电荷的电场线：电场线从正电荷出发，终止于负电荷}\label{fig:dipole-field-lines}
\end{center}

\subsection*{模型 2：等量同种电荷的电场}

设两个等量正电荷位于 $x$ 轴上的 $(-a,0)$、$(a,0)$，中垂线上一点
$P(0,y)$ 到两电荷的距离均为 $r=\sqrt{a^2+y^2}$。两电荷在 $P$ 点产生
的场强水平方向分量相互抵消，竖直方向分量相加，因此当 $y>0$ 时
\[
E(y)=2\left(k\frac{Q}{r^2}\right)\frac{y}{r}
=\frac{2kQy}{(a^2+y^2)^{3/2}},
\]
方向沿中垂线向上；$y<0$ 时方向相反。由表达式或函数图像可知：中点
$y=0$ 处场强为零，离开中点后先增大，在 $|y|=a/\sqrt2$ 处达到最大，
随后逐渐减小并趋近于零。

\begin{center}
% TikZ 图：等量同种电荷的电场线
% Day 15：电荷、电场与电场强度
\begin{tikzpicture}[scale=0.9,>=Latex]
  \coordinate (A) at (-1.4,0);
  \coordinate (B) at (1.4,0);

  \filldraw[red] (A) circle (4pt);
  \node[below,fill=white,inner sep=1pt] at (A) {$+Q$};
  \filldraw[red] (B) circle (4pt);
  \node[below,fill=white,inner sep=1pt] at (B) {$+Q$};

  % 同种正电荷的电场线都从电荷出发，向远处延伸。
  % 面向彼此的一侧会向上下分流；电场线不连接、不相交。
  \draw[->,thick] (-1.58,0.00) -- (-3.10,0.00);
  \draw[->,thick] (-1.53,0.15) .. controls (-1.95,0.35) and (-2.45,0.75) .. (-3.05,1.25);
  \draw[->,thick] (-1.53,-0.15) .. controls (-1.95,-0.35) and (-2.45,-0.75) .. (-3.05,-1.25);
  \draw[->,thick] (-1.30,0.18) .. controls (-1.05,0.55) and (-0.88,1.15) .. (-0.82,1.95);
  \draw[->,thick] (-1.30,-0.18) .. controls (-1.05,-0.55) and (-0.88,-1.15) .. (-0.82,-1.95);

  \draw[->,thick] (1.58,0.00) -- (3.10,0.00);
  \draw[->,thick] (1.53,0.15) .. controls (1.95,0.35) and (2.45,0.75) .. (3.05,1.25);
  \draw[->,thick] (1.53,-0.15) .. controls (1.95,-0.35) and (2.45,-0.75) .. (3.05,-1.25);
  \draw[->,thick] (1.30,0.18) .. controls (1.05,0.55) and (0.88,1.15) .. (0.82,1.95);
  \draw[->,thick] (1.30,-0.18) .. controls (1.05,-0.55) and (0.88,-1.15) .. (0.82,-1.95);

  \filldraw[black] (0,0) circle (1.6pt);
  \node[below,fill=white,inner sep=1pt] at (0,0) {$O$};
  \node[font=\small] at (0,-2.20) {中点处两场强等大反向，合场强为零};
\end{tikzpicture}


\figcaption{等量同种正电荷的电场线：中点场强为零}\label{fig:equal-like-field-lines}

\vspace{0.4em}
\begin{tikzpicture}[x=1.25cm,y=1.15cm,>=Stealth,line cap=round,line join=round]
  \draw[->,thick] (-0.15,0) -- (3.65,0) node[right] {$|y|/a$};
  \draw[->,thick] (0,-0.12) -- (0,1.35) node[above] {$E$};

  \draw[very thick,red!80!black,samples=120,domain=0:3.35]
    plot (\x,{2.60*\x/(1+\x*\x)^(1.5)});

  \pgfmathsetmacro{\xm}{1/sqrt(2)}
  \draw[dashed,gray] (\xm,0) -- (\xm,1.0);
  \draw[dashed,gray] (0,1.0) -- (\xm,1.0);
  \fill[red!80!black] (\xm,1.0) circle (1.8pt);

  \node[below,font=\scriptsize] at (\xm,0) {$1/\sqrt2$};
  \node[left,font=\scriptsize] at (0,1.0) {$E_{\max}$};
  \node[font=\scriptsize,align=center] at (2.55,0.86)
    {$E=\dfrac{2kQ|y|}{(a^2+y^2)^{3/2}}$};
  \node[font=\scriptsize,gray!75!black] at (2.55,0.22)
    {$|y|\to\infty$ 时，$E\to0$};
\end{tikzpicture}


\figcaption{等量同种正电荷的中垂线场强：从零增大到最大值后再减小}\label{fig:equal-like-bisector-field}
\end{center}

\subsection*{模型 3：匀强电场}

平行板电容器之间的电场近似为匀强电场（两极板间电场线平行等距）。
本节先根据平行、等距的电场线认识匀强电场，并直接使用题目给出的
$E$。电势差与场强的定量关系将在 Day 16 建立。

\begin{center}
% TikZ 图：平行板间的匀强电场线
% Day 15：电荷、电场与电场强度
\begin{tikzpicture}[scale=0.9,>=Latex]
  \draw[very thick] (-2.8,1.0) -- (2.8,1.0);
  \draw[very thick] (-2.8,-1.0) -- (2.8,-1.0);
  \node[left] at (-2.8,1.0) {$+$};
  \node[left] at (-2.8,-1.0) {$-$};

  \foreach \x in {-2,-1,0,1,2} {
    \draw[->,thick,blue] (\x,0.78) -- (\x,-0.78);
  }
  \draw[<->] (3.2,-1.0) -- (3.2,1.0);
  \node[right] at (3.2,0) {$d$};
  \node[font=\small] at (0,-1.45) {板间中部电场线近似平行等距，$U$ 为两板间电压};
\end{tikzpicture}


\figcaption{平行板间的匀强电场线：中部区域近似平行等距}\label{fig:parallel-plate-field-lines}
\end{center}

\subsection*{模型 4：静电平衡与静电屏蔽}

导体中有可以自由移动的电荷。把导体放入外电场后，自由电荷会重新分布，直到导体内部的合电场为零，这种状态叫\textbf{静电平衡}。

\begin{center}
\begin{tikzpicture}[scale=0.92,>=Stealth,line cap=round,line join=round]
  \path[use as bounding box] (-4.35,-2.05) rectangle (4.35,2.15);

  % 绕过导体的外部场线。
  \draw[blue!70!black,->,thick] (-4.1,1.70) -- (4.1,1.70);
  \draw[blue!70!black,->,thick] (-4.1,-1.70) -- (4.1,-1.70);
  \draw[blue!70!black,->,thick]
    (-4.1,1.03) .. controls (-2.75,1.02) and (-2.10,1.50) ..
    (0,1.56) .. controls (2.10,1.50) and (2.75,1.02) .. (4.1,1.03);
  \draw[blue!70!black,->,thick]
    (-4.1,-1.03) .. controls (-2.75,-1.02) and (-2.10,-1.50) ..
    (0,-1.56) .. controls (2.10,-1.50) and (2.75,-1.02) .. (4.1,-1.03);

  % 一部分场线终止于左侧感应负电荷，另一部分从右侧正电荷出发。
  \foreach \y/\xe in {-0.56/1.43,0/1.65,0.56/1.43}{
    \draw[blue!70!black,->,thick] (-4.1,\y) -- (-\xe,\y);
    \draw[blue!70!black,->,thick] (\xe,\y) -- (4.1,\y);
  }
  \node[blue!70!black,font=\small] at (-3.25,2.02)
    {外电场方向};

  % 导体壳与空腔。
  \fill[gray!22] (0,0) ellipse (1.65 and 1.12);
  \draw[thick] (0,0) ellipse (1.65 and 1.12);
  \fill[white] (0,0) ellipse (0.92 and 0.56);
  \draw[thick] (0,0) ellipse (0.92 and 0.56);

  \foreach \y in {-0.57,-0.19,0.19,0.57}{
    \node[blue!75!black,font=\small] at (-1.48,\y) {$-$};
    \node[red!80!black,font=\small] at (1.48,\y) {$+$};
  }

  \node[font=\small] at (0,1.29) {导体壳};
  \node[font=\small,align=center] at (0,0) {空腔内无电荷\\$E=0$};

  \draw[->,gray!75!black,thick] (-2.20,-1.88) -- (-1.18,-0.82);
  \node[font=\scriptsize,align=center] at (-2.70,-1.90)
    {感应电荷\\重新分布};
  \draw[->,gray!75!black,thick] (2.20,-1.88) -- (1.10,-0.78);
  \node[font=\scriptsize,align=center] at (2.72,-1.90)
    {导体材料内部\\$E=0$};
\end{tikzpicture}


\figcaption{静电屏蔽：导体静电平衡时内部电场为零；空腔内无电荷时，可隔离外电场影响}\label{fig:electrostatic-shielding}
\end{center}

高中阶段需要理解三条结论：
\begin{enumerate}
  \item 静电平衡时，导体内部电场强度为零。
  \item 导体上多余的净电荷分布在导体表面；若有空腔，还要区分内、外表面。
  \item 导体表面附近的电场方向垂直于表面；否则电荷还会沿表面移动，说明尚未平衡。
\end{enumerate}

\textbf{学习边界：} 本节重点是判断导体内部场强、表面电荷分布趋势和屏蔽效果，不要求计算复杂导体表面的电荷密度分布。看到静电屏蔽题，先判断导体是否已经达到静电平衡；若讨论空腔内的场强，还要先确认空腔内是否放有电荷。无内部电荷时，封闭导体空腔可隔离外电场；若腔内有电荷，腔内一般仍有电场。

\textbf{应用理解：} 金属外壳、屏蔽线、汽车车厢和法拉第笼都利用静电屏蔽保护内部区域。这里不是“外面没有电场”，而是导体表面感应电荷产生的附加电场抵消了导体内部的外电场。

\textbf{尖端放电与防雷装置：} 导体静电平衡时，净电荷分布在表面；在尖锐位置，电荷更容易集中，附近电场更强。当电场强到足以使空气电离时，空气由绝缘状态变得能导电，就会发生尖端放电或电晕。建筑物防雷系统通常由接闪器、引下线和接地装置组成，其核心作用是在雷击发生时优先承接雷电流，并把雷电流经低阻通路安全导入大地。尖端附近可能发生电晕，但不能把防雷作用简单解释成“靠尖端放电消除云层电荷”。

\begin{center}
% TikZ 图：尖端放电与避雷针
% Day 15：电荷、电场与电场强度
\begin{tikzpicture}[scale=0.95,>=Latex,line cap=round,line join=round]
  \path[use as bounding box] (-4.4,-1.65) rectangle (4.4,2.4);

  % Charged cloud.
  \draw[thick,gray!70] (-3.9,1.85)
    .. controls (-3.6,2.25) and (-3.0,2.20) .. (-2.75,1.88)
    .. controls (-2.35,2.25) and (-1.72,2.10) .. (-1.58,1.76)
    .. controls (-1.05,1.95) and (-0.62,1.66) .. (-0.78,1.32)
    .. controls (-1.55,1.18) and (-2.75,1.20) .. (-3.9,1.35)
    .. controls (-4.15,1.48) and (-4.16,1.70) .. (-3.9,1.85);
  \foreach \x/\y in {-3.35/1.65,-2.70/1.78,-2.05/1.55,-1.35/1.60} {
    \node[blue!70!black,font=\small] at (\x,\y) {$-$};
  }
  \node[font=\small] at (-2.35,2.28) {带电云层};

  % Ground and building.
  \draw[thick] (-4.2,-0.95) -- (4.2,-0.95);
  \foreach \x in {-4.0,-3.55,...,4.0} {
    \draw[gray!60] (\x,-0.95) -- ++(-0.18,-0.12);
  }
  \draw[fill=gray!12,thick] (1.3,-0.95) rectangle (2.75,0.25);
  \draw[fill=gray!18,thick] (1.18,0.25) -- (2.02,0.82) -- (2.88,0.25) -- cycle;
  \node[font=\small] at (2.03,-0.35) {建筑};

  % Lightning rod.
  \draw[very thick] (2.02,0.82) -- (2.02,1.65);
  \draw[fill=black] (2.02,1.65) -- (1.92,1.42) -- (2.12,1.42) -- cycle;
  \draw[very thick] (2.02,0.82) -- (2.02,-1.20);
  \node[font=\scriptsize,right=2pt] at (2.05,-1.18) {接地};

  % Dense positive charges near sharp tip.
  \foreach \x/\y in {1.86/1.47,1.95/1.56,2.05/1.56,2.17/1.47,2.02/1.33} {
    \node[red!80!black,font=\scriptsize] at (\x,\y) {$+$};
  }

  % Field lines around the sharp point.
  \foreach \a in {72,92,112,132} {
    \draw[red!75!black,->,thick] (2.02,1.55) -- ++(\a:0.90);
  }
  \node[red!80!black,font=\small,align=center] at (3.45,1.88)
    {尖端附近\\电场更强};

  % Smooth conductor comparison.
  \draw[fill=gray!18,thick] (-2.10,-0.95) rectangle (-1.25,0.28);
  \draw[fill=gray!18,thick] (-2.10,0.28) arc[start angle=180,end angle=0,radius=0.425] -- cycle;
  \foreach \x/\y in {-1.98/0.48,-1.72/0.62,-1.45/0.48} {
    \node[red!80!black,font=\scriptsize] at (\x,\y) {$+$};
  }
  \draw[red!75!black,->] (-1.95,0.55) -- (-2.35,0.95);
  \draw[red!75!black,->] (-1.68,0.68) -- (-1.68,1.15);
  \draw[red!75!black,->] (-1.42,0.55) -- (-1.02,0.95);
  \node[font=\small,align=center] at (-1.68,-1.35)
    {圆滑表面\\电荷较分散};

\end{tikzpicture}


\figcaption{尖端电场与建筑防雷：接闪器、引下线和接地装置构成安全泄流通路}\label{fig:point-discharge-lightning-rod}
\end{center}

\subsection*{模型 5：匀强电场中的等效重力场}

若带电物体同时受到重力和匀强电场力，且两个力在整个运动区域内大小、方向都不变，就可以把它们先合成为一个恒力。对质量为 $m$、电荷量为 $q$ 的物体：
\[
\vect{F}_{\text{合}}=m\vect{g}+q\vect{E}=m\vect{g'} .
\]
这里的 $\vect{g'}$ 叫等效重力场强度，它不是新的真实重力，而是“重力和电场力合起来的效果”。若 $\vect{E}$ 水平向右且 $q>0$，则
\[
g'=\sqrt{g^2+\left(\frac{qE}{m}\right)^2},\qquad
\tan\theta=\frac{qE}{mg}.
\]

\begin{center}
% TikZ 图：匀强电场中的等效重力场
% Day 15：电荷、电场与电场强度
\begin{tikzpicture}[scale=0.86,>=Latex,line cap=round,line join=round,
  every node/.style={font=\small}]
  % Left: force composition on a positive charge.
  \begin{scope}[shift={(-3.1,0)}]
    \draw[blue!55,->] (-1.45,1.25) -- (1.45,1.25) node[right] {$\vect{E}$};
    \foreach \y in {0.62,-0.62} {
      \draw[blue!55,->] (-1.45,\y) -- (1.45,\y);
    }
    \filldraw[red] (0,0) circle (3pt);
    \node[red,above left=4pt] at (0,0) {$+q$};
    \draw[very thick,red!80!black,->] (0,0) -- (1.15,0);
    \node[red!80!black,above=4pt] at (0.78,0) {$qE$};
    \draw[very thick,black,->] (0,0) -- (0,-1.15) node[left] {$mg$};
    \draw[very thick,orange!85!black,->] (0,0) -- (1.15,-1.15)
      node[below right] {$mg'$};
    \draw[dashed,gray] (1.15,0) -- (1.15,-1.15) -- (0,-1.15);
  \end{scope}

  % Right: pendulum aligns with the resultant force.
  \begin{scope}[shift={(2.8,0)}]
    \coordinate (O) at (0,1.35);
    \coordinate (P) at (0.82,-0.65);
    \filldraw[black] (O) circle (1.4pt);
    \draw[dashed,gray] (O) -- (0,-1.15);
    \draw[thick] (O) -- (P);
    \filldraw[red] (P) circle (4pt);
    \node[red,above left=4pt] at (P) {$+q$};
    \draw[blue!55,->] (-1.45,0.88) -- (1.45,0.88) node[right] {$\vect{E}$};
    \draw[blue!55,->] (-1.45,0.25) -- (1.45,0.25);
    \draw[very thick,orange!85!black,->] (P) -- ++(0.60,-1.10)
      node[right] {$mg'$};
    \draw[very thick,black,->] (P) -- ++(0,-0.95) node[left] {$mg$};
    \draw[very thick,red!80!black,->] (P) -- ++(0.90,0);
    \node[red!80!black,above=4pt] at ($(P)+(0.66,0)$) {$qE$};
    \draw[->] (0,0.66) arc[start angle=-90,end angle=-68,radius=0.58];
    \node at (0.30,0.42) {$\theta$};
  \end{scope}
\end{tikzpicture}


\figcaption{等效重力场：把重力和电场力先合成，再按力学模型处理}\label{fig:equivalent-gravity-field}
\end{center}

\textbf{模型价值：} 带电小球悬挂在匀强电场中、带电小球沿斜面运动、复合场中的平衡问题，都可以先判断电场力方向，再与重力合成。平衡时，绳的拉力、支持力等约束力用来平衡这个合力。

\textbf{模型边界：} 只有电场近似匀强、重力近似恒定时，才能把合力看成固定方向的等效重力。若电场来自点电荷，方向和大小随位置变化，就不能简单用一个固定的 $g'$ 代替。

\subsection*{模型 6：带电液滴的悬浮}

宏观带电小球、带电液滴和尘埃颗粒在电场中运动时，重力通常不能忽略。若它在平行板间保持静止或匀速直线运动，就说明合力为零。最常见的悬浮条件是
\[
|q|E=mg.
\]
\begin{center}
% TikZ 图：带电液滴在匀强电场中悬浮
% Day 15：电荷、电场与电场强度
\begin{tikzpicture}[scale=0.88,>=Latex,line cap=round,line join=round,
  every node/.style={font=\small}]
  \draw[thick] (-2.1,1.2) -- (2.1,1.2);
  \draw[thick] (-2.1,-1.2) -- (2.1,-1.2);
  \node[left] at (-2.1,1.2) {$+$};
  \node[left] at (-2.1,-1.2) {$-$};

  \foreach \x in {-1.55,-0.75,0.05,0.85,1.65} {
    \draw[blue!70,->] (\x,0.85) -- (\x,-0.85);
  }
  \node[blue!70] at (1.95,0.12) {$\vect{E}$};

  \filldraw[orange!85!black] (0,0) circle (4pt);
  \node[orange!85!black,left=5pt] at (0,0) {$-q$};
  \draw[very thick,red!80!black,->] (0,0) -- (0,0.88);
  \node[red!80!black,right=5pt] at (0,0.65) {$F_{\text{电}}=|q|E$};
  \draw[very thick,black,->] (0,0) -- (0,-0.88);
  \node[black,right=5pt] at (0,-0.65) {$mg$};

  \draw[gray,<->] (2.42,-1.2) -- (2.42,1.2);
  \node[gray,right] at (2.5,0) {$d$};

  \node[align=center] at (0,-1.72)
    {悬浮条件：$F_{\text{电}}=mg$，即 $|q|E=mg$};
\end{tikzpicture}


\figcaption{带电液滴悬浮：电场力与重力平衡，电荷正负决定电场力方向}\label{fig:charged-droplet-suspension}
\end{center}

\textbf{认知提醒：} 先由电场方向和电荷正负判断电场力方向，再列平衡方程。若图中电场向下，负电荷所受电场力向上；正电荷则受力向下，不能悬浮在这个方向的电场中。

\section*{六、例题精讲}

\begin{examplebox}{库仑定律的应用}

\textbf{题意分析：} 两个点电荷 $Q_1 = 2 \times 10^{-8}\ \mathrm{C}$、$Q_2 = -3 \times 10^{-8}\ \mathrm{C}$ 相距 $r = 0.1\ \mathrm{m}$。求它们之间的库仑力大小和方向。

\textbf{物理模型：} 库仑定律。

\begin{center}
\begin{tikzpicture}[x=1cm,y=1cm,>=Stealth,line cap=round,line join=round]
  \coordinate (A) at (-2.15,0);
  \coordinate (B) at (2.15,0);

  \fill[red!80!black] (A) circle (0.075);
  \fill[blue!75!black] (B) circle (0.075);
  \node[above=5pt,red!80!black] at (A) {$Q_1>0$};
  \node[above=5pt,blue!75!black] at (B) {$Q_2<0$};

  \draw[->,very thick,red!80!black] ($(A)+(0.18,0)$) -- ($(A)+(1.15,0)$);
  \node[below=3pt,red!80!black] at (-1.47,0) {$\vect F_{12}$};
  \draw[->,very thick,blue!75!black] ($(B)+(-0.18,0)$) -- ($(B)+(-1.15,0)$);
  \node[below=3pt,blue!75!black] at (1.47,0) {$\vect F_{21}$};

  \draw[<->,gray,thick] (-2.15,-0.72) -- (2.15,-0.72)
    node[midway,below] {$r$};
  \node[font=\scriptsize] at (0,0.58) {异种电荷相互吸引};
\end{tikzpicture}


\figcaption{例题图：异种点电荷相互吸引，两个库仑力等大反向}\label{fig:example-coulomb-attraction}
\end{center}

\textbf{解题过程：}

\[
F = k\frac{|Q_1 Q_2|}{r^2} = 9 \times 10^9 \times \frac{2 \times 10^{-8} \times 3 \times 10^{-8}}{(0.1)^2}
\]
\[
= 9 \times 10^9 \times \frac{6 \times 10^{-16}}{0.01} = 9 \times 10^9 \times 6 \times 10^{-14} = 5.4 \times 10^{-4}\ \mathrm{N}
\]

方向：异种电荷相互吸引。

\textbf{方法总结：} 库仑力题先用绝对值计算大小，再由电荷正负判断方向；最后检查指数运算和单位。

\end{examplebox}

\begin{examplebox}{电场强度的计算}

\textbf{题意分析：} 在 $x$ 轴上 $x_1 = -0.1\ \mathrm{m}$ 处放 $Q_1 = +4 \times 10^{-9}\ \mathrm{C}$，$x_2 = 0.1\ \mathrm{m}$ 处放 $Q_2 = -4 \times 10^{-9}\ \mathrm{C}$。求原点 $O$ 处的电场强度。

\textbf{物理模型：} 等量异种电荷，电场叠加。

\begin{center}
\begin{tikzpicture}[x=1cm,y=1cm,>=Stealth,line cap=round,line join=round]
  \draw[->,thick] (-3.0,0) -- (3.25,0) node[right] {$x$};
  \coordinate (A) at (-2.0,0);
  \coordinate (O) at (0,0);
  \coordinate (B) at (2.0,0);

  \fill[red!80!black] (A) circle (0.075);
  \fill[black] (O) circle (0.05);
  \fill[blue!75!black] (B) circle (0.075);
  \node[below=5pt,red!80!black] at (A) {$Q_1>0$};
  \node[below=5pt] at (O) {$O$};
  \node[below=5pt,blue!75!black] at (B) {$Q_2<0$};

  \draw[->,very thick,red!80!black] (0,0.34) -- (1.05,0.34)
    node[midway,above] {$\vect E_1$};
  \draw[->,very thick,blue!75!black] (0,-0.34) -- (1.05,-0.34)
    node[midway,below] {$\vect E_2$};
  \draw[->,ultra thick,black] (0,0) -- (1.75,0)
    node[near end,above] {$\vect E$};

  \draw[<->,gray] (-2.0,0.83) -- (0,0.83) node[midway,above] {$0.10\ \mathrm m$};
  \draw[<->,gray] (0,0.83) -- (2.0,0.83) node[midway,above] {$0.10\ \mathrm m$};
\end{tikzpicture}


\figcaption{例题图：原点处两个分场强方向相同}\label{fig:example-dipole-origin-field}
\end{center}

\textbf{解题过程：}

$Q_1$ 在 $O$ 点产生场强：方向向右（$+Q \to$ 正电荷场强方向向外）。
\[
E_1 = k\frac{Q_1}{r^2} = 9 \times 10^9 \times \frac{4 \times 10^{-9}}{(0.1)^2} = 9 \times 10^9 \times 4 \times 10^{-7} = 3600\ \mathrm{N/C}
\]

$Q_2$ 在 $O$ 点产生场强：方向向右（$-Q \to$ 负电荷场强方向向内）。
\[
E_2 = k\frac{|Q_2|}{r^2} = 3600\ \mathrm{N/C}\ (\text{方向向右})
\]

合场强：$E = E_1 + E_2 = 7200\ \mathrm{N/C}$，方向向右。

\textbf{方法总结：} 计算合场强时，先分别计算各电荷产生的场强大小和方向，再用矢量加法求和。在同一直线上用正负表示方向。

\end{examplebox}

\begin{examplebox}{电场线与运动轨迹}

\textbf{题意分析：} 一个正电荷仅在电场力作用下从静止释放，它会沿电场线运动吗？

\textbf{解题过程：}
不一定。正电荷从静止释放的瞬间，加速度方向沿该点电场线的切线方向。
若这条电场线是直线，且此后各处场强始终沿这条直线，电荷会沿直线运动；
若电场线是曲线，电场力只决定每一时刻的加速度方向，速度还保留此前运动的
结果，因此粒子一般不会沿曲线电场线前进。

\textbf{拓展说明：} 若电荷具有非零初速度，还要同时考虑初速度方向；即使电场线
是直线，横向初速度也会使轨迹偏离电场线。

\textbf{方法总结：} 看到电场线先把它当“场强方向示意图”，不当轨迹。电场线给出
加速度的瞬时方向，粒子轨迹由速度和加速度共同决定。

\end{examplebox}

\begin{examplebox}{匀强电场中的带电小球平衡}

\textbf{题意分析：} 质量为 $m$、电荷量为 $+q$ 的小球用绝缘细线悬挂在
水平向右的匀强电场中。静止时细线与竖直方向的夹角为 $\theta$，求
$\theta$ 和细线拉力 $T$。

\begin{center}
\begin{tikzpicture}[x=1cm,y=1cm,>=Stealth,line cap=round,line join=round]
  \coordinate (O) at (0,1.55);
  \coordinate (P) at (0.86,-0.62);

  \foreach \y in {1.05,0.35,-0.35,-1.05}{
    \draw[blue!55,->] (-2.15,\y) -- (2.35,\y);
  }
  \node[blue!70!black] at (2.55,1.05) {$\vect E$};

  \fill[black] (O) circle (1.5pt);
  \draw[dashed,gray] (O) -- (0,-1.22);
  \draw[thick] (O) -- (P);
  \fill[red!80!black] (P) circle (0.09);
  \node[red!80!black,left=5pt] at (P) {$+q$};

  \draw[->,very thick,orange!85!black] (P) -- ($(P)!0.42!(O)$);
  \node[orange!85!black,above left=2pt] at ($(P)!0.25!(O)$) {$\vect T$};
  \draw[->,very thick,red!80!black] (P) -- ++(1.18,0);
  \node[red!80!black,below right=2pt] at ($(P)+(1.18,0)$) {$q\vect E$};
  \draw[->,very thick,black] (P) -- ++(0,-1.02)
    node[left=3pt] {$m\vect g$};

  \draw[->] (0,0.78) arc[start angle=-90,end angle=-68.4,radius=0.78];
  \node at (0.25,0.64) {$\theta$};
\end{tikzpicture}


\figcaption{例题图：带电小球在重力、电场力和拉力作用下平衡}\label{fig:example-charged-pendulum-field}
\end{center}

\textbf{解题过程：} 对小球受力分析，水平方向和竖直方向分别有
\[
T\sin\theta=qE,\qquad T\cos\theta=mg.
\]
两式相除并平方相加，得到
\[
\boxed{\tan\theta=\frac{qE}{mg}},\qquad
\boxed{T=\sqrt{(mg)^2+(qE)^2}}.
\]

\textbf{方法总结：} 先由电荷正负判断电场力方向，再建立水平、竖直坐标轴。
若电荷改为负电，偏转方向反向，但拉力大小仍由两个恒力的合力大小决定。

\end{examplebox}

\section*{七、变式训练}

\begin{enumerate}
  \item 两个点电荷 $Q_1 = 3 \times 10^{-8}\ \mathrm{C}$、$Q_2 = 5 \times 10^{-8}\ \mathrm{C}$ 相距 $0.2\ \mathrm{m}$，求库仑力大小。

  \item 在 $x$ 轴上 $x = 0.2\ \mathrm{m}$ 处有一个 $Q = -2 \times 10^{-9}\ \mathrm{C}$ 的点电荷。求 $x = 0$ 处的电场强度大小和方向。
  \begin{center}
  \begin{tikzpicture}[scale=0.95,>=Stealth,line cap=round,line join=round]
  \draw[->,thick] (-0.55,0) -- (4.55,0) node[right] {$x$};
  \coordinate (A) at (0,0);
  \coordinate (Q) at (3.0,0);
  \draw[thick] (A) ++(0,0.08) -- ++(0,-0.16);
  \draw[thick] (Q) ++(0,0.08) -- ++(0,-0.16);
  \draw[fill=red!70,draw=red!80!black] (Q) circle (0.075);
  \node[below,font=\small,align=center] at (A) {场点\\$x=0$};
  \node[below,font=\small,align=center] at (Q) {$x=0.20\,\mathrm m$};
  \node[above,red,font=\small] at (Q) {$Q=-2\times10^{-9}\,\mathrm C$};
  \draw[<->,gray,thick] (0,0.82) -- (3,0.82)
    node[midway,above,font=\small] {$r=0.20\,\mathrm m$};
  \draw[->,blue,thick] (0,0.32) -- (1.15,0.32)
    node[midway,above] {$\vec E$};
  \node[font=\small] at (1.55,1.36) {负电荷产生的场强指向电荷};
\end{tikzpicture}

  \end{center}

  \item 等量同种正电荷 $+Q$ 相距 $2a$，求它们连线中点处的电场强度。
  \begin{center}
  % TikZ 图：等量同种电荷的电场线
% Day 15：电荷、电场与电场强度
\begin{tikzpicture}[scale=0.9,>=Latex]
  \coordinate (A) at (-1.4,0);
  \coordinate (B) at (1.4,0);

  \filldraw[red] (A) circle (4pt);
  \node[below,fill=white,inner sep=1pt] at (A) {$+Q$};
  \filldraw[red] (B) circle (4pt);
  \node[below,fill=white,inner sep=1pt] at (B) {$+Q$};

  % 同种正电荷的电场线都从电荷出发，向远处延伸。
  % 面向彼此的一侧会向上下分流；电场线不连接、不相交。
  \draw[->,thick] (-1.58,0.00) -- (-3.10,0.00);
  \draw[->,thick] (-1.53,0.15) .. controls (-1.95,0.35) and (-2.45,0.75) .. (-3.05,1.25);
  \draw[->,thick] (-1.53,-0.15) .. controls (-1.95,-0.35) and (-2.45,-0.75) .. (-3.05,-1.25);
  \draw[->,thick] (-1.30,0.18) .. controls (-1.05,0.55) and (-0.88,1.15) .. (-0.82,1.95);
  \draw[->,thick] (-1.30,-0.18) .. controls (-1.05,-0.55) and (-0.88,-1.15) .. (-0.82,-1.95);

  \draw[->,thick] (1.58,0.00) -- (3.10,0.00);
  \draw[->,thick] (1.53,0.15) .. controls (1.95,0.35) and (2.45,0.75) .. (3.05,1.25);
  \draw[->,thick] (1.53,-0.15) .. controls (1.95,-0.35) and (2.45,-0.75) .. (3.05,-1.25);
  \draw[->,thick] (1.30,0.18) .. controls (1.05,0.55) and (0.88,1.15) .. (0.82,1.95);
  \draw[->,thick] (1.30,-0.18) .. controls (1.05,-0.55) and (0.88,-1.15) .. (0.82,-1.95);

  \filldraw[black] (0,0) circle (1.6pt);
  \node[below,fill=white,inner sep=1pt] at (0,0) {$O$};
  \node[font=\small] at (0,-2.20) {中点处两场强等大反向，合场强为零};
\end{tikzpicture}

  \end{center}

  \item 质量为 $m$、电荷量为 $+q$ 的小球用细线悬挂在水平向右的匀强电场中，静止时细线与竖直方向成角 $\theta$。写出 $\tan\theta$ 与 $E$、$m$、$q$、$g$ 的关系。

  \item 一封闭导体壳处于静电平衡。判断下列说法：\\
  （1）外电场中，导体材料内部的场强是否为零？\\
  （2）空腔内无电荷时，外电场能否在空腔内产生静电场？\\
  （3）若把一点电荷放入空腔，空腔内场强是否仍处处为零？
\end{enumerate}

\begin{learningbox}{变式训练参考结果}
\begin{enumerate}[itemsep=0.45em,topsep=0.3em,parsep=0pt]
  \item $F=k|Q_1Q_2|/r^2=3.38\times10^{-4}\ \mathrm N$。
  \item $E=k|Q|/r^2=450\ \mathrm{N/C}$，方向沿 $x$ 轴正方向，即指向负电荷。
  \item 两个正电荷在中点产生的电场强度大小相等、方向相反，合场强为零。
  \item 平衡时 $T\sin\theta=qE$、$T\cos\theta=mg$，故
  $\tan\theta=qE/(mg)$。
  \item （1）导体材料内部场强为零；（2）不能，封闭空腔可隔离外电场；
  （3）不能，腔内有电荷时空腔内一般存在电场，但导体材料内部仍为零。
\end{enumerate}
\end{learningbox}

\section*{八、竞赛入门补充}

\begin{competitionbox}
\textbf{竞赛视角 1：静电力作用下的平衡问题}

库仑力可以与绳的拉力、重力等共同平衡。竞赛入门题常把“力的分解”和“几何关系”放在一起考查。

\textbf{例题：} 两个质量均为 $m$、带电量均为 $+q$ 的小球，用等长细线悬挂在同一点。平衡时两线夹角为 $2\theta$，求 $q$。

\begin{center}
% TikZ 图：两个带电小球的静电平衡
% Day 15：电荷、电场与电场强度
\begin{tikzpicture}[scale=0.78,>=Latex]
  \coordinate (O) at (0,1.8);
  \coordinate (L) at (-1.45,-0.45);
  \coordinate (R) at (1.45,-0.45);

  \filldraw[black] (O) circle (1.5pt);
  \draw[thick] (O) -- (L) node[midway,left] {$l$};
  \draw[thick] (O) -- (R) node[midway,right] {$l$};

  \draw[dashed,gray] (O) -- (0,-0.75);
  \draw[->] (0,1.25) arc[start angle=-90,end angle=-57,radius=0.55];
  \node at (0.32,1.05) {$\theta$};

  \filldraw[red] (L) circle (4pt);
  \filldraw[red] (R) circle (4pt);
  \node[below] at (L) {$+q$};
  \node[below] at (R) {$+q$};

  \draw[<->] (-1.45,-0.95) -- (1.45,-0.95);
  \node[below] at (0,-0.95) {$r$};

  \draw[->,thick,blue] (R) -- (2.25,-0.45) node[right] {$F$};
  \draw[->,thick,blue] (R) -- (1.02,0.22) node[right] {$T$};
  \draw[->,thick,blue] (R) -- (1.45,-1.08) node[right] {$mg$};
\end{tikzpicture}


\figcaption{带电小球平衡：库仑斥力、重力与绳拉力共同平衡}\label{fig:charged-pendulum-equilibrium}
\end{center}

\textbf{解：} 对一个小球受力分析：
\begin{itemize}
  \item 水平方向：库仑斥力 $F = k\dfrac{q^2}{r^2}$，绳拉力的水平分力 $T\sin\theta$。
  \item 竖直方向：重力 $mg$，绳拉力的竖直分力 $T\cos\theta$。
\end{itemize}
由平衡条件：
\[
T\sin\theta = k\frac{q^2}{r^2},\qquad T\cos\theta = mg
\]
两式相除，并利用两球间距 $r=2l\sin\theta$（$l$ 为绳长），得到
\[
\tan\theta=\frac{kq^2}{mgr^2},\qquad
\boxed{q=2l\sin\theta\sqrt{\frac{mg\tan\theta}{k}}}.
\]

\medskip
\hrule
\medskip
\textbf{竞赛视角 2：用对称性快速求合场强}

\textbf{例题：} 正方形边长为 $a$，四个顶点原来各放一个相同正点电荷 $+q$。若去掉右上角的电荷，求正方形中心 $O$ 处的合场强大小和方向。

\begin{center}
% TikZ 图：正方形缺角电荷的对称补全
% Day 15：电荷、电场与电场强度
\begin{tikzpicture}[scale=0.78,>=Latex]
  \coordinate (A) at (-1.25,1.25);
  \coordinate (B) at (1.25,1.25);
  \coordinate (C) at (1.25,-1.25);
  \coordinate (D) at (-1.25,-1.25);
  \coordinate (O) at (0,0);

  \draw[thick] (A) -- (B) -- (C) -- (D) -- cycle;
  \filldraw[red] (A) circle (3.5pt) node[above left] {$+q$};
  \draw[red,thick] (B) circle (5pt);
  \node[red,above right] at (B) {缺};
  \filldraw[red] (C) circle (3.5pt) node[below right] {$+q$};
  \filldraw[red] (D) circle (3.5pt) node[below left] {$+q$};

  \filldraw[black] (O) circle (2pt);
  \node[below left] at (O) {$O$};
  \draw[dashed,gray] (O) -- (B);
  \draw[->,very thick,blue] (O) -- (0.9,0.9) node[midway,above left] {$\vect{E}$};
  \draw[<->] (-1.25,-1.65) -- (1.25,-1.65);
  \node[below] at (0,-1.65) {$a$};
\end{tikzpicture}


\figcaption{正方形缺角模型：用“完整对称系统”反推剩余电荷的合场强}\label{fig:square-missing-charge-field}
\end{center}

\textbf{解：} 四个顶点都放 $+q$ 时，中心处的合场强为零。去掉右上角电荷后，剩下三个电荷的合场强等于“原来被去掉的那个电荷在 $O$ 处场强”的反向。中心到顶点的距离为 $r=a/\sqrt2$，因此
\[
E_0=k\frac{q}{r^2}
=k\frac{q}{(a/\sqrt2)^2}
=\frac{2kq}{a^2}.
\]
所以剩下三个电荷在中心处的合场强大小也是
$E=2kq/a^2$，方向指向被去掉电荷所在的顶点。
\end{competitionbox}

\section*{九、易错点}

\begin{mistakebox}
\begin{enumerate}
  \item \textbf{库仑定律中 $r$ 和 $r^2$ 混用：} 先量两点电荷间距离 $r$，再把这个距离平方后放入分母；不要把 $r$ 本身直接当作分母。

  \item \textbf{误以为摩擦能创造电荷：} 分析起电现象时先看电子从哪里转移到哪里；起电过程只是电子转移或电荷重新分布，总电荷守恒。

  \item \textbf{混淆 $E = F/q_0$ 和 $E = k|Q|/r^2$：} 前者是定义式（用正试探电荷定义任意电场），后者是点电荷的场强大小公式（只适用于点电荷），方向还要单独判断。

  \item \textbf{用试探电荷 $q$ 的正负判断场强方向：} 先按正试探电荷受力方向确定场强方向；若放入负电荷，它受力方向与场强方向相反。

  \item \textbf{电场叠加时忘了 $E$ 的方向：} 先分别画出每个场强方向，再用平行四边形定则或分量法叠加。

  \item \textbf{认为电场线是电荷的轨迹：} 先区分“场的方向图”和“粒子运动轨迹”。电场线只是描述工具，不是运动轨迹。

  \item \textbf{误以为静电屏蔽会消灭外电场：} 静电平衡时，导体材料内部的
  合电场为零；封闭空腔只有在腔内无电荷时，才能隔离外电场。屏蔽并不会
  使导体外部的电场消失。

  \item \textbf{把避雷针理解成“吸引雷电”或“消除雷电”：} 防雷系统由接闪器、
  引下线和接地装置组成，主要作用是在雷击发生时提供可控的低阻泄流通路。
  尖端电晕不是消除云层电荷、保证不发生雷击的手段。

  \item \textbf{等效重力场中忘记电荷正负：} $q>0$ 时电场力沿电场方向，$q<0$ 时电场力反向。先判力的方向，再合成。
\end{enumerate}
\end{mistakebox}

\section*{十、今日小结}

\begin{enumerate}
  \item 电荷有正、负两种；摩擦起电、接触起电和感应起电都遵循总电荷守恒。
  \item 库仑力大小 $F = k\dfrac{|Q_1 Q_2|}{r^2}$，
  $k = 9.0\times10^9\ \mathrm{N\cdot m^2/C^2}$。
  \item 电场强度 $\vect E = \dfrac{\vect F}{q_0}$（定义式，$q_0$ 为正试探电荷；大小为 $E=\dfrac{F}{q_0}$），点电荷场强大小 $E = k\dfrac{|Q|}{r^2}$。
  \item 电场强度是矢量，遵循叠加原理。
  \item 电场线：从正电荷或无穷远出发，终止于负电荷或无穷远；同一示意图中疏密定性表示强弱，切线表示方向。
  \item 匀强电场的场强大小和方向处处相同，电场线平行等距。
  \item 导体静电平衡时，导体材料内部电场为零；空腔内无电荷时可隔离外电场。
  尖端附近电场较强，可能使空气发生电离。
  \item 匀强电场中的带电物体可把 $mg$ 与 $qE$ 先合成为等效重力；带电液滴悬浮时常用 $|q|E=mg$。
\end{enumerate}

\textbf{一分钟复盘：} 看到电场题，先在题旁写三项：场源是谁、研究的是场强还是电场力、试探电荷正负；再画出场强叠加方向或等效重力方向。方向确认后再选择公式，避免把“场的方向”和“电荷受力方向”混在一起。
